\section{Introduction}
\label{sec:intro}

Retrieval Augmented Generation (RAG) emerges as a promising paradigm for minimizing hallucinations in Large Language Models (LLMs) by based on external knowledge~\cite{lewis2020retrieval,shuster2021retrieval,ji2023survey,gao2023retrieval,asai2023self,qiu2025efficient,tao2024llms,linghu2025llm}. While standard RAG systems excel at retrieving explicit information, they often struggle with complex queries that require multi-hop reasoning across disparate documents. To address this, Graph-based RAG (GraphRAG) approaches~\cite{edge2024local} have been introduced, which organize the corpus as a structured graph to facilitate multi-step information propagation.

Despite their effectiveness, existing GraphRAG methods face a significant efficiency bottleneck~\cite{han2025rag,peng2024graph}. Traditional approaches typically rely on LLMs to extract entity-relation triplets and then construct the knowledge graphs. This process incurs prohibitive computational costs, which makes them impractical for large-scale corpus. To mitigate this, recent lightweight frameworks such as LinearRAG~\cite{zhuang2025linearrag} have proposed replacing expensive LLM-based relation extraction with efficient Named Entity Recognition (NER) and modeling document structures directly.

However, we argue that these lightweight approaches introduce a critical limitation: the \textbf{Semantic Gap}. By relying solely on explicit structural co-occurrence (i.e., entities appearing in the same sentence or document), these methods fail to capture latent connections between semantically related but structurally disjoint entities. Consider the example illustrated in Figure~\ref{fig:teaser}. To answer the question \textit{'Who is the spouse of the current monarch of the UK?'}, a system must connect 'monarch' in document B to 'Queen' in document A. Structure-only graph would treat these as distinct nodes because they are not same phrase and never appear in the same context window. Consequently, the reasoning chain is broken, which leads to retrieval failure. 

To bridge this gap without sacrificing efficiency, we propose \textbf{\EHGRAG} (\underline{E}fficient \underline{H}ypergraph-based \underline{RAG}), a novel framework that unifies structural and semantic indexing via hypergraphs. Unlike simple graphs, hypergraphs use hyperedges to connect arbitrary sets of nodes, making them naturally suitable for modeling both explicit document inclusion (structural hyperedges) and implicit semantic clusters (semantic hyperedges). Specifically, we employ lightweight NER for node extraction and construct structural hyperedges based on sentence-level co-occurrence. Simultaneously, we cluster entity embeddings to form semantic hyperedges, which connect semantically similar entities across different documents. During retrieval, we introduce a hybrid retrieval mechanism that utilizes diffusion-based entity activation to propagate query influence through explicit and latent pathways, followed by topic-aware passage scoring and PPR refinement to accurately identify the top-$k$ relevant documents based on both content and graph structure.

\begin{figure}[t]
    \centering
    \includegraphics[width=\columnwidth]{figure/Figure1.pdf} 
    \caption{An example of the semantic gap. While structural graph fails to link disjoint entities (monarch and Queen), \EHGRAG bridges them via a semantic hyperedge, enabling the retrieval of the multi-hop reasoning.}
    \label{fig:teaser}
\vspace{-3mm}
\end{figure}

To validate the effectiveness of \EHGRAG, we conduct experiments on four multi-hop QA benchmarks. Experimental results demonstrate that our method significantly outperforms state-of-the-art baselines including HippoRAG2 and LinearRAG, improving accuracy by up to 6.9\% on 2WikiMultiHop while maintaining zero token consumption for indexing.

In summary, our contributions are as follows:
\begin{itemize}
\item We identify the \textit{semantic gap} in existing lightweight GraphRAG methods, where the lack of semantic connectivity compromises multi-hop reasoning.
\item We propose \EHGRAG, a hybrid hypergraph-based RAG framework that efficiently integrates structural co-occurrence and latent semantic correlations via entity activation by hypergraph diffusion and well designed passage scoring strategy.
\item Extensive experiments across four benchmarks demonstrate that \EHGRAG achieves state-of-the-art performance with linear indexing complexity, offering a scalable solution for knowledge-intensive tasks.
\end{itemize}